\section{Conclusions}\label{sec:conclu}

Our politico-economic equilibrium model showcases the impact of various pension system features on the size of transfers and, directly and indirectly, savings and labor supply. In particular, the model takes into account two key aspects of pension design. First, we considered systems that either have or do not have a requirement for years of contributions to receive full benefits. 
Second, we considered whether benefits are earnings-related or flat, thus allowing for both Bismarckian and Beveridgean systems. 
 
With logarithmic GHH preferences the model features an analytical characterization of equilibrium taxes as a function only of demographics and parameters. We find, in line with previous work, that equilibrium taxes increase with the political power of retirees, population ageing, and income inequality. Taxes also increase when pensions are Beveridgean and in more universal systems, but only when the share of informal households benefiting from the latter is large enough. 


Calibrating the model to Argentina, it accounts for almost half of the observed increase in pension spending due to a reform that made pensions universal, and for more over time as the newly covered households retire with less private savings. The model predicts a decrease in private savings of 0.3 p.p.\ of GDP and a decrease in labor supply of about 0.3\%, in line with observed evidence. Taxation is the main factor behind the reduction in private savings and labor supply. 

A calibration to the United States, France, and the United Kingdom suggests that ageing and the degree with which benefits are linked to contributions are the main determinants of social security taxes, with income (and voting) inequality playing a minor role. Differences in preferences for leisure, that in the model capture pure preferences and differences in labor regulations and institutions, play a minor role for equilibrium taxes. 

An important limitation of our framework is the absence of risk, even though uncertainty is a key driver of both savings and labor supply decisions over the life cycle. In related work, \textcite{BergGEwp}, we analyze a similar environment in which households are homogeneous when young but face health risks that may prevent some from qualifying for contributive pensions at retirement. By allowing the probability of pension eligibility to depend on labor supply, we show that households respond to this risk by working and saving more than they otherwise would.\footnote{When evaluating the pension reform in Argentina, we find, in that context, a higher response of savings and labor supply. Conversely, given the lack of heterogeneity, the effects of changes in the degree by which benefits are tied to earnings has negligible equilibrium effects.} We regard the present analysis and that earlier work as complementary in shedding light on different channels through which social security systems shape behavior, and how their design affects its politico-economic equilibrium. 


Our model could be valuable in exploring the challenges posed by the expected rise of new types of platform work, such as zero-hour contracts with no guaranteed hours, which make formal labor less attractive.
This raises concerns about the pension coverage for workers engaged in these activities, as these workers face low job security under existing legislation.

Finally, we let the electorate choose the design of the system as well as its size. Costless redistribution within a cohort is politically too attractive for an interior design to survive; once redistributive funds carry a deadweight cost calibrated to the observed U.S.\ system, ageing pushes the chosen design toward earnings-related benefits, a flatter income distribution moves it to the Bismarckian corner, and a flatter distribution of voting participation pulls it the other way. Which of the last two wins depends on the intertemporal elasticity of substitution, so the model can reproduce or reverse the weak cross-country relation between inequality and pension design, and the data we use cannot tell the two apart. 



